%===============================================================================
\section{Introduction}
\label{sec:introduction}
%===============================================================================

Understanding how cells respond to genetic perturbations, such as those induced by CRISPR screens~\citep{dixit2016perturb}, is fundamental to advancing our understanding of basic biology and therapeutic discovery~\citep{plenge2013validating,adamson2016multiplexed,norman2019exploring}. 
However, the combinatorial explosion across $\sim$20,000 protein-coding genes makes exhaustive perturbation experiments infeasible~\citep{norman2019exploring}, motivating computational models that predict cellular responses to unseen perturbations or combinations thereof.

Existing approaches to predict the transcriptional outcome following a perturbation face a fundamental challenge: generative quality comes at the cost of causal interpretability.
For example, recent methods like GEARS~\citep{roohani2022gears} or scGPT~\citep{cui2024scgpt} achieve strong predictive performance in cellular response to perturbations; however, they act as ``black boxes'' providing little mechanistic insight into how perturbations affect cellular systems.
This opacity prevents researchers from understanding why certain combinations of perturbations produce synergistic effects, failing to provide testable biological hypotheses.


As such, computational biology demands models that provide mechanistic insights into cellular function~\citep{chen_applying_2024}, making \gls{crl}~\citep{scholkopf2021toward} a promising approach for mapping genes to biological processes.
In particular, recent \gls{crl} approaches such as \gls{dvae}~\citep{zhang2023identifiability} or SENA~\citep{de2025interpretable} aim at recovering the underlying causal generative factors and their interactions, enabling counterfactual reasoning.
However, these \gls{crl} approaches suffer from fundamental architectural limitations that prevent them from fulfilling their theoretical causal identifiability guarantees: the use of a dense interventional encoder inherently means that every gene, including those intervened upon can influence all latent variables.
We term this phenomenon \emph{intervention spillover}, as it renders the task of identifying a unique latent target for an intervention ill-posed.
Consequently, any mechanism designed to apply the intervention is forced to converge to a suboptimal solution, because it cannot perform the clean, single-node manipulation required for valid causal inference~\citep{brehmer2022weakly,lippe2022citris}.

\textcolor{blue}{Further, the destructive nature of \gls{scRNAseq} makes it impossible to observe the same cell before and after a genetic intervention, complicating the prediction of cell-specific responses.
Existing methods address this issue by randomly pairing control and perturbed cells.
Since the model must minimize error across all possible arbitrary pairings, it is incentivized to learn the population-mean response.
We refer to this phenomenon as \emph{mean collapse}, as this regression to the mean erodes cellular diversity.}

We present \gls{raptorgraph}, a novel framework that addresses causal identifiability challenges in perturbation prediction through structured encoder design. 
Our first key innovation is the GraphPathway layer, a novel architecture that uses preconditioning to enforce sparse mappings from genes to learned Causal Meta-Pathways. 
This design enables the model to perform the clean, single-node manipulations required for valid causal inference, directly maintaining interpretability.
This innovation is complemented by a novel \gls{ot} preprocessing step, which mitigates the distinct problem of mean collapse.
We benchmark the generative quality of \glsentryshort{raptorgraph} on the \citet{norman2019exploring} Perturb-seq dataset \textcolor{blue}{and the \citet{replogle_combinatorial_2020} CRISPRi dataset.
For the \citet{norman2019exploring} data, we further evaluate the} accuracy of genetic interaction in perturbation combinations, \textcolor{blue}{and the biological interpretability of the inferred Causal Meta-Pathway.
Specifically, when applied to scRNA-seq data from a K562 cell line, \glsentryshort{raptorgraph} uncovered distinct, biologically accurate programs, such as the differentiation-induced cell cycle arrest driven by EGR1 and stress-induced growth arrest mediated by JUN.}
\textcolor{blue}{Additionally, we \citet{replogle_combinatorial_2020} CRISPRi dataset}
We demonstrate strong performance alongside interpretability, which addresses key gaps not filled by current approaches.
